\documentclass{beamer}
% COMSC-3013 Computer Architecture shared Beamer preamble.
% Week decks live one directory below weeks/ and input this file as:
%   % COMSC-3013 Computer Architecture shared Beamer preamble.
% Week decks live one directory below weeks/ and input this file as:
%   % COMSC-3013 Computer Architecture shared Beamer preamble.
% Week decks live one directory below weeks/ and input this file as:
%   \input{../_shared/beamer-preamble.tex}
% Keep the course self-contained. Change visual/build conventions here,
% not by forking one-off week themes.

\usetheme{Boadilla}
\usecolortheme{seahorse}
\setbeamertemplate{navigation symbols}{}
\setbeamertemplate{footline}[frame number]

\usepackage{amsmath}
\usepackage{booktabs}
\usepackage{graphicx}
\usepackage{listings}
\usepackage{pgfpages}

\lstset{
  basicstyle=\ttfamily\small,
  columns=fullflexible,
  keepspaces=true,
  showstringspaces=false,
  breaklines=true
}

% Speaker notes remain hidden in the ordinary student build. Define \shownotes
% before inputting the deck to create an instructor copy; see weeks/README.md.
\ifdefined\shownotes
  \setbeameroption{show notes on second screen=right}
\fi

\newcommand{\ArchTitleSlide}[4]{%
  % #1 week label, #2 title, #3 central question, #4 Wednesday prediction
  \begin{frame}
    \titlepage
  \end{frame}
  \begin{frame}{The machine question}
    \Large #3

    \vspace{1.2em}
    \normalsize\textbf{Prediction to carry into Wednesday:} #4
  \end{frame}
}

% Keep the course self-contained. Change visual/build conventions here,
% not by forking one-off week themes.

\usetheme{Boadilla}
\usecolortheme{seahorse}
\setbeamertemplate{navigation symbols}{}
\setbeamertemplate{footline}[frame number]

\usepackage{amsmath}
\usepackage{booktabs}
\usepackage{graphicx}
\usepackage{listings}
\usepackage{pgfpages}

\lstset{
  basicstyle=\ttfamily\small,
  columns=fullflexible,
  keepspaces=true,
  showstringspaces=false,
  breaklines=true
}

% Speaker notes remain hidden in the ordinary student build. Define \shownotes
% before inputting the deck to create an instructor copy; see weeks/README.md.
\ifdefined\shownotes
  \setbeameroption{show notes on second screen=right}
\fi

\newcommand{\ArchTitleSlide}[4]{%
  % #1 week label, #2 title, #3 central question, #4 Wednesday prediction
  \begin{frame}
    \titlepage
  \end{frame}
  \begin{frame}{The machine question}
    \Large #3

    \vspace{1.2em}
    \normalsize\textbf{Prediction to carry into Wednesday:} #4
  \end{frame}
}

% Keep the course self-contained. Change visual/build conventions here,
% not by forking one-off week themes.

\usetheme{Boadilla}
\usecolortheme{seahorse}
\setbeamertemplate{navigation symbols}{}
\setbeamertemplate{footline}[frame number]

\usepackage{amsmath}
\usepackage{booktabs}
\usepackage{graphicx}
\usepackage{listings}
\usepackage{pgfpages}

\lstset{
  basicstyle=\ttfamily\small,
  columns=fullflexible,
  keepspaces=true,
  showstringspaces=false,
  breaklines=true
}

% Speaker notes remain hidden in the ordinary student build. Define \shownotes
% before inputting the deck to create an instructor copy; see weeks/README.md.
\ifdefined\shownotes
  \setbeameroption{show notes on second screen=right}
\fi

\newcommand{\ArchTitleSlide}[4]{%
  % #1 week label, #2 title, #3 central question, #4 Wednesday prediction
  \begin{frame}
    \titlepage
  \end{frame}
  \begin{frame}{The machine question}
    \Large #3

    \vspace{1.2em}
    \normalsize\textbf{Prediction to carry into Wednesday:} #4
  \end{frame}
}


\title{Week 4 Wednesday}
\subtitle{Use the operating system as a telescope}
\author{COMSC-3013 Computer Architecture}
\date{}

\begin{document}

\begin{frame}
  \titlepage
\end{frame}

\ArchQuestionSlide{What is this computer doing?}{A small set of observations can answer bounded subquestions, but no one command reveals the whole physical machine.}
\note{Open with the vague question. Ask students what they would need to know before trusting an answer.}

\begin{frame}{The useful model: question $\rightarrow$ instrument $\rightarrow$ boundary}
\begin{itemize}
  \item Decompose one vague machine question into observable subquestions.
  \item Choose an instrument whose output bears on exactly one subquestion.
  \item Interpret the observation inside its execution scope; name what it cannot prove.
\end{itemize}
\note{This is Lens 4 in action: decomposition makes evidence and limits visible.}
\end{frame}

\begin{frame}[fragile]{Start with the course Observatory}
\begin{lstlisting}
./lab/bin/archprobe --out-dir /tmp/arch-week04
cat /tmp/arch-week04/machine.txt
\end{lstlisting}
\begin{itemize}
  \item The Observatory labels the visible execution environment.
  \item It gives us a common receipt before platform-specific instruments.
\end{itemize}
\note{Run archprobe first in the recording. Scope is part of the evidence, not an afterthought.}
\end{frame}

\begin{frame}{One vague question, six smaller questions}
\begin{itemize}
  \item Which OS and kernel are visible? \hfill \texttt{uname}
  \item Which ISA / machine label is visible? \hfill \texttt{uname -m}, \texttt{lscpu}
  \item How many logical processors are visible? \hfill \texttt{lscpu}
  \item How much memory is visible? \hfill \texttt{free}
  \item What processes are running? \hfill \texttt{ps}
  \item What kind of file or binary is this? \hfill \texttt{file}
\end{itemize}
\note{Do not turn this into a command vocabulary test. Each command earns its place by answering a question.}
\end{frame}

\begin{frame}[fragile]{Persistent worked example: one file, two views}
\begin{lstlisting}
printf 'ABC\n' > /tmp/arch-week04-bytes.txt
file /tmp/arch-week04-bytes.txt
od -An -tx1 -c /tmp/arch-week04-bytes.txt
\end{lstlisting}
\begin{itemize}
  \item \textbf{Observation:} a tool reports a text file and bytes.
  \item \textbf{Interpretation:} the bytes are consistent with the visible text specimen.
  \item \textbf{Boundary:} this does not explain every file, encoding, or running process.
\end{itemize}
\note{Keep the same specimen visible while moving from file identity to raw representation.}
\end{frame}

\begin{frame}{Tempting claim: one command tells the whole story}
\Large
\textbf{Wrong:} ``\texttt{lscpu} says x86\_64, so I know the physical computer.''

\vspace{1em}
\normalsize
\textbf{Why it fails:} the output describes the interface exposed to this execution environment. Containers, virtual machines, permissions, and runtime policy can hide or reshape details.
\note{Make the tempting shortcut explicit, then point back to the scope label from archprobe.}
\end{frame}

\begin{frame}[fragile]{Add instruments selectively}
\begin{lstlisting}
uname -a
lscpu
free -h
ps -eo pid,comm,%cpu,%mem --sort=-%cpu | head
file /bin/sh
\end{lstlisting}
\begin{itemize}
  \item Run only what helps answer the chosen subquestion.
  \item If an instrument is unavailable, use the Observatory/fallback packet and record the limitation.
\end{itemize}
\note{The goal is an evidence chain, not a shell-command catalog.}
\end{frame}

\begin{frame}{The Explain / Defend receipt}
\begin{enumerate}
  \item Question: what were you trying to know?
  \item Instrument: which tool produced the observation?
  \item Observation: quote only the evidence that matters.
  \item Interpretation: what does it support?
  \item Boundary: what does it not establish?
  \item Revision: how would you ask more precisely next time?
\end{enumerate}
\note{Friday asks students to defend one strong claim, not submit a screenshot dump.}
\end{frame}

\begin{frame}{Carry one bounded claim into Friday}
\Large
Choose one subquestion. Capture one raw observation. Explain one interpretation. Name one boundary.

\vspace{1em}
\normalsize
The machine can answer a precise question; it cannot rescue an imprecise one.
\note{End with the receipt structure. Leave the strongest claim for students to make and defend.}
\end{frame}

\end{document}
